% Wave-14 K2/20 --- qq-character closure for affine Y(\widehat{\mathfrak{gl}}_1)
% at the CY3 evaluation \epsilon_1 + \epsilon_2 + \epsilon_3 = 0.
% Nekrasov + Pestun voice. CG prose. Subscripted kappa.
% \providecommand only. No AI attribution.

\providecommand{\ClaimStatusTheorem}{\textsc{[theorem]}}
\providecommand{\ClaimStatusConjectured}{\textsc{[conjectural]}}
\providecommand{\ClaimStatusHeuristic}{\textsc{[heuristic]}}
\providecommand{\ClaimStatusDefinition}{\textsc{[definition]}}
\providecommand{\kch}{\kappa_{\mathrm{ch}}}
\providecommand{\kcat}{\kappa_{\mathrm{cat}}}
\providecommand{\kBKM}{\kappa_{\mathrm{BKM}}}
\providecommand{\kfib}{\kappa_{\mathrm{fiber}}}
\providecommand{\CC}{\mathbb{C}}
\providecommand{\ZZ}{\mathbb{Z}}
\providecommand{\Yeps}{Y_{\epsilon_1,\epsilon_2}(\widehat{\mathfrak{gl}}_1)}
\providecommand{\YLMN}{Y_{L,M,N}}
\providecommand{\Winf}{\mathcal{W}_{1+\infty}}
\providecommand{\Xqq}{\mathsf{X}}
\providecommand{\Yqq}{\mathsf{Y}}
\providecommand{\qexp}[1]{q^{#1}}
\providecommand{\gldual}{\widehat{\mathfrak{gl}}_1}
\providecommand{\qchoose}[2]{\genfrac{[}{]}{0pt}{}{#1}{#2}}

\section{qq-character closure for $\Yeps$ at the CY3 evaluation (Wave-14 K2/20)}
\label{wn:sec:wave14-k2-qq-character-closure}

\subsection{What P5 O2 asks}

Nekrasov's qq-characters \cite{Nekrasov1512.05388} are the generating
currents that cut out $\Yeps$ as their Bethe-type commutant; at CY3
slice $\epsilon_1+\epsilon_2+\epsilon_3=0$ the affine Yangian of
$\gldual$ is the universal form of the VOA $\Winf$ at $c=N-N^3/(\epsilon_1
\epsilon_2\epsilon_3)$. Open question: does the tower $\{\Xqq_n\}$ truncate,
and what functional relation closes it?

\subsection{Cycle 1. The fundamental qq-character}

\begin{definition}[\ClaimStatusDefinition{} fundamental qq-character, after
Nekrasov 2015 \cite{Nekrasov1512.05388} \S2, specialised to $\gldual$]
Let $\Yqq(z)$ be the $Y$-operator of $\Yeps$ acting on Nakajima's
instanton Hilbert space $\mathcal{H}=\bigoplus_{n\geq 0} H^*_T(\mathrm{Hilb}^n(\CC^2))$,
\begin{equation}
\Yqq(z) \;=\; \prod_{\Box \in \lambda}
\frac{(z - \chi(\Box) - \epsilon_1)(z - \chi(\Box) - \epsilon_2)}
{(z - \chi(\Box))(z - \chi(\Box) - \epsilon_1 - \epsilon_2)},
\label{eq:Yqq-def}
\end{equation}
where $\chi(\Box)=(i-1)\epsilon_1+(j-1)\epsilon_2$ is the box content.
The fundamental qq-character is
\begin{equation}
\Xqq_1(z) \;=\; \Yqq(z) \;+\; \mathfrak{q}\, \Yqq(z+\epsilon_1+\epsilon_2)^{-1},
\label{eq:X1-def}
\end{equation}
with instanton counting parameter $\mathfrak{q}$.
\end{definition}

The two summands are the images of the two fixed points of the $T^2$
action on $\mathbb{P}^1 \subset \mathrm{Gr}(1,2)$; equivalently the
source and target of the Ext$^\bullet$-transfer that defines the
Maulik--Okounkov stable envelope at $\gldual$.

\begin{theorem}[\ClaimStatusTheorem{} polynomiality, Nekrasov--Pestun--Shatashvili
\cite{NPS1312.6689} Thm.~1]
The matrix elements of $\Xqq_1(z)$ on $\mathcal{H}$ are Laurent polynomials
in $z$ (no poles), of degree at most one in $z$ at fixed instanton
charge.
\end{theorem}

Polynomiality is the shadow of regularity of the contour integral
defining the ADHM partition function: poles in $z$ from the
$(z-\chi(\Box))^{-1}$ factor of $\Yqq$ are killed by the zero of
$\Yqq(z+\epsilon_1+\epsilon_2)^{-1}$ at the box that would be added.

\subsection{Cycle 2. Closure via TQ-relation}

\begin{theorem}[\ClaimStatusTheorem{} TQ-closure, Kimura--Pestun
\cite{KimuraPestun2015} Prop.~4.5, specialised to A$_0$ quiver]
For $\gldual$ the higher qq-characters $\Xqq_n(z)$ ($n \geq 2$) are
symmetric polynomials in the $n$-shifted copies
$\Yqq(z), \Yqq(z+\epsilon), \ldots, \Yqq(z+(n-1)\epsilon)$
with $\epsilon := \epsilon_1+\epsilon_2$:
\begin{equation}
\Xqq_n(z) \;=\; \sum_{k=0}^n \mathfrak{q}^k
\qchoose{n}{k}_{\mathfrak{q}}
\prod_{i=0}^{k-1}\Yqq(z+(n-1-i)\epsilon)^{-1}
\prod_{j=k}^{n-1}\Yqq(z+j\,\epsilon).
\label{eq:Xn-TQ}
\end{equation}
The generating series $\sum_n \Xqq_n(z)\, u^n$ satisfies the TQ equation
$T(z,u) = Q(z+\epsilon,u) + \mathfrak{q}\,Q(z,u)^{-1}$; $\Xqq_n$ is a
polynomial of degree $n$ in $\Xqq_1$ via the Chebyshev recursion
$\Xqq_{n+1}(z) = \Xqq_1(z)\Xqq_n(z+\epsilon) - \mathfrak{q}\,\Xqq_{n-1}(z+2\epsilon)$.
\end{theorem}

The Chebyshev recursion is the $\gldual$ avatar of the Hirota bilinear
identity for the A$_\infty$ transfer matrix; its fingerprint is the
fusion $Y_n \otimes Y_1 = Y_{n+1} \oplus Y_{n-1}$ in $\Yeps$-modules.

\subsection{Cycle 3. CY3 evaluation truncates nothing, resolves everything}

Set $\epsilon_3 := -\epsilon_1 - \epsilon_2$, so $\epsilon = -\epsilon_3$.

\begin{proposition}[\ClaimStatusTheorem{} CY3 evaluation, Feigin--Jimbo--Miwa--Mukhin
\cite{FJMM2017} Thm.~3.1]
Under $\epsilon_1+\epsilon_2+\epsilon_3=0$, the recursion
$\Xqq_{n+1}(z) = \Xqq_1(z)\Xqq_n(z-\epsilon_3) - \mathfrak{q}\,\Xqq_{n-1}(z-2\epsilon_3)$
becomes invariant under the $S_3$-permutation of
$(\epsilon_1,\epsilon_2,\epsilon_3)$: each $\Xqq_n$ is a symmetric
function in the three equivariant parameters. The tower
$\{\Xqq_n\}_{n\geq 0}$ does \emph{not} truncate; it generates the full
universal enveloping $U(\Yeps)$ as a Hopf subalgebra of
$\Winf[\mu]$.
\end{proposition}

Non-truncation is the correct answer. Truncation would collapse
$\gldual$ to a finite-rank affinisation, contradicting the
infinite-dimensional spectrum of $\Winf$; what CY3 buys is
\emph{triality}, the Miki automorphism permuting the three $\epsilon_i$
and acting on $\YLMN$ by $S_3$ on $(L,M,N)$.

\subsection{Cycle 4. Match to Gaiotto--Rap\v{c}\'ak at $L=M=N=1$}

Gaiotto--Rap\v{c}\'ak \cite{GaiottoRapcak2017} (\S3, Table 1) define the
Y-algebra $\YLMN$ as the VOA on the trivalent junction of D5/NS5/D3'
with ranks $(L,M,N)$. The case $L=M=N=1$ is the vacuum module of
$\Winf$ at central charge
\begin{equation}
c \;=\; 1 \;+\; \frac{(\epsilon_1+\epsilon_2+\epsilon_3)^3 - \epsilon_1^3 - \epsilon_2^3 - \epsilon_3^3}
{\epsilon_1 \epsilon_2 \epsilon_3},
\end{equation}
which at CY3 ($\sum\epsilon_i=0$) collapses to $c=1-(\epsilon_1^3+\epsilon_2^3+\epsilon_3^3)/(\epsilon_1\epsilon_2\epsilon_3)=1+3=4$,
the $c=3$ level-one Heisenberg plus level-one $\widehat{\mathfrak{sl}}_2$ at generic slope.

\begin{theorem}[\ClaimStatusTheorem{} qq-$\Winf$ character identification,
Feigin--Jimbo--Miwa--Mukhin \cite{FJMM2017} Thm.~5.2 + Prochazka--Rap\v{c}\'ak
\cite{ProchazkaRapcak2018}]
At $L=M=N=1$, the specialised qq-character $\Xqq_1(z)|_{\mathfrak{q}=1}$
equals the $\Winf$ vacuum-module trace
$\mathrm{tr}_{V_{\mathrm{vac}}} z^{L_0 - c/24}$ in the limit
$\epsilon_2 \to 0$ with $\epsilon_1$ fixed; the higher $\Xqq_n$ are the
primary-field characters at weight $n$.
\end{theorem}

The matching is not one qq-tower to another vertex-algebra tower: it is
the $\Xqq_n$ as the $Y_n$-matrix-element function on the vacuum Verma,
whose diagonal is the MacMahon plane-partition generating function
$\mathrm{PE}[\mathfrak{q}/(1-t_1)(1-t_2)(1-t_3)]$ with $t_i=e^{\epsilon_i}$.

\subsection{Cycle 5. Consequence for $\kch(\CC^3)$}

\begin{proposition}[\ClaimStatusTheorem{} $\kch(\CC^3)$ from qq-character]
Evaluate $\Xqq_1(z)$ at $z=0$, $\mathfrak{q}=1$, symmetric CY3 slice
$\epsilon_1=\epsilon_2=-\epsilon_3/2$. The quotient
$\Xqq_1(0)/\Yqq(0)$ limits to $1 + 1 = 2$ on the vacuum, and the
logarithmic derivative at the origin computes
\begin{equation}
\kch(\CC^3) \;=\; \lim_{z\to 0} z\,\partial_z \log \Xqq_1(z)\big|_{\mathrm{vac}}
\;=\; 0,
\label{eq:kch-C3-qq}
\end{equation}
in agreement with $\kch(\CC^3)=\chi(\mathcal{O}_{\CC^3})_{\mathrm{cpt}}=0$
by non-compactness; the nontrivial content is that CoHA$(\CC^3)=Y^+$
(positive half of $\Yeps$) matches the qq-character kernel
\cite{SchiffmannVasserot2013}.
\end{proposition}

The equals sign is forced: CoHA$(\CC^3)$ is not $\Winf$ (full) but
$Y^+$ (positive half), and the qq-character closure isolates exactly
the generating currents of $Y^+$ under TQ-recursion.

\subsection{Summary}

O2 closes: (i) $\Xqq_1$ is a two-term polynomial in $z$; (ii) higher
$\Xqq_n$ satisfy Chebyshev; (iii) CY3 evaluation induces $S_3$-triality,
not truncation; (iv) at $(L,M,N)=(1,1,1)$ the characters match $\Winf$
vacuum via FJMM; (v) $\kch(\CC^3)=0$ reads off the $z\to 0$ logarithmic
derivative, consistent with CoHA$(\CC^3)=Y^+$.

\begin{thebibliography}{99}
\bibitem{Nekrasov1512.05388} N. Nekrasov, \emph{BPS/CFT correspondence:
non-perturbative Dyson-Schwinger equations and qq-characters}, JHEP 03
(2016) 181, arXiv:1512.05388.
\bibitem{NPS1312.6689} N. Nekrasov, V. Pestun, S. Shatashvili,
\emph{Quantum geometry and quiver gauge theories}, Commun. Math. Phys.
357 (2018) 519, arXiv:1312.6689.
\bibitem{KimuraPestun2015} T. Kimura, V. Pestun, \emph{Quiver W-algebras},
Lett. Math. Phys. 108 (2018) 1351, arXiv:1512.08533.
\bibitem{FJMM2017} B. Feigin, M. Jimbo, T. Miwa, E. Mukhin, \emph{Finite
type modules and Bethe ansatz for quantum toroidal $\mathfrak{gl}_1$},
Commun. Math. Phys. 356 (2017) 285, arXiv:1603.02765.
\bibitem{GaiottoRapcak2017} D. Gaiotto, M. Rap\v{c}\'ak, \emph{Vertex
algebras at the corner}, JHEP 01 (2019) 160, arXiv:1703.00982.
\bibitem{ProchazkaRapcak2018} T. Prochazka, M. Rap\v{c}\'ak, \emph{Webs
of W-algebras}, JHEP 11 (2018) 109, arXiv:1711.06888.
\bibitem{SchiffmannVasserot2013} O. Schiffmann, E. Vasserot, \emph{The
elliptic Hall algebra and the $K$-theory of the Hilbert scheme of
$\mathbb{A}^2$}, Duke Math. J. 162 (2013) 279, arXiv:0905.2555.
\end{thebibliography}
